\section{语音大模型与实时口语对话}
\label{sec:nlp-speech-realtime-dialogue}

前面几节的识别、理解、合成与翻译模型可以串成一个口语对话系统，但这样的系统有三处缺口。其一，理解能力局限于各部件的训练任务，转写之外的语气、情绪和环境声在文字接口处丢失。其二，识别、文本生成、合成串联，延迟逐段相加。其三，系统按话轮交替工作，无法边听边说，被打断后也要靠外部检测才停下。\term{全双工}（full duplex）指双方可以同时发送和接收声音的交互条件，它要求模型在说话时仍持续读取用户的声音。

方法谱系沿两条线推进：声音怎样进入语言模型，以及声音怎样从模型产生。表\ref{tab:nlp-speech-dialogue-routes}列出三类路线。级联系统用识别器把声音变成文字，交给文本对话模型，再用合成器读出；\term{半级联}（half-cascade）系统把语音编码器直接接到大语言模型上，声音特征进入模型，输出仍是文字，再经外部合成；全双工联合模型把用户与系统两路声音都离散成词元，用同一个模型在共同时间轴上持续建模。半级联解决声音怎样进入模型；级联与半级联仍要靠外部的话轮检测与响应门控决定何时说，这两部分在全双工模型之前讨论；无论哪条路线，声音回复都要与业务动作协调。

\begin{table*}[htbp]
  \centering
  \begin{tabularx}{\linewidth}{@{}>{\raggedright\arraybackslash}p{20mm}>{\raggedright\arraybackslash}p{30mm}>{\raggedright\arraybackslash}X>{\raggedright\arraybackslash}X>{\raggedright\arraybackslash}X>{\raggedright\arraybackslash}X@{}}
    \toprule
    路线 & 组成 & 延迟来源 & 打断处理 & 副语言信息 & 可检查性\\
    \midrule
    级联 & 识别器—文本对话模型—合成器 & 等识别提交、文本生成、合成首段，三段相加 & 依赖外部VAD与轮次检测，停播由控制逻辑执行 & 在转写处丢失，合成端另设语气 & 转写与回复文字都可检查\\
    半级联 & 语音编码器—适配器—大语言模型，外接合成 & 编码整段声音后生成文字，再加合成 & 同级联，仍需外部检测 & 输入端保留，输出端仍由合成器决定 & 回复文字可检查，输入只能看模型表现\\
    全双工联合 & 两路声音词元与文字流，单一模型逐帧生成 & 每步一个编解码帧加声学延迟 & 用户声音持续输入，模型自行停止或交叠 & 输入与输出都在声音词元中 & 文字流可读；工具与业务状态需外接\\
    \bottomrule
  \end{tabularx}
  \caption{口语对话的三类路线。从上到下，声音在模型内部保留得越来越多，延迟中的串联部分越来越少，但对话行为也越来越依赖模型自身，越难用显式规则控制。表中延迟一列只列结构来源，不是实测数值。}
  \label{tab:nlp-speech-dialogue-routes}
\end{table*}

\subsection{语音编码器与语言模型的连接}
\label{sec:nlp-speech-llm-connector}

大语言模型已经具备指令理解与推理能力，转写却会丢掉语气、情绪与非语音事件。半级联路线的思路是不经过文字，直接把声音的连续表示接到语言模型的输入嵌入空间。困难在于两种表示的长度与分布都不匹配：Whisper编码器（第\ref{sec:nlp-speech-whisper}小节）每20\,ms输出一帧，30\,s录音有1500帧，而同样长度的转写只有一两百个词元；编码器的特征空间也与语言模型的词嵌入空间无关。

连接结构由三部分组成。预训练语音编码器把声音编码为$\bm H\in\mathbb R^{T\times d_a}$；适配器（adapter，指编码器与语言模型之间的连接模块，不同于第\ref{sec:nlp-sft-peft}节插在层内的适配器）把帧率降到每秒十余个到数十个向量，并投影到语言模型的宽度$d_{\mathrm{LM}}$；得到的音频嵌入与文字指令的词嵌入拼成一条序列送入语言模型。常见的降采样方式有两类。第一类是拼接或池化：把相邻$k$帧拼成一个$k d_a$维向量，再用两层MLP投影，
\begin{equation}
 \bm e_m=\mathrm{MLP}\bigl([\bm h_{(m-1)k+1};\ldots;\bm h_{mk}]\bigr)\in\mathbb R^{d_{\mathrm{LM}}},\qquad m=1,\ldots,\lceil T/k\rceil.
 \label{eq:nlp-speech-adapter-stack}
\end{equation}
第二类是查询读取：Q-Former用一组可学习查询向量，以交叉注意力从编码器特征中读出固定数目的输出向量。整段只用一组查询会丢失时间顺序，语音中常改为窗口级：把$\bm H$切成长$L$帧的窗口，每个窗口用$N$个查询读出$N$个向量，总长度为$N\lceil T/L\rceil$，保持单调的时间顺序。视觉—语言模型中的Q-Former与投影连接器采用同一框架：冻结的感知编码器后接一个小型可训练连接模块，把可变长特征变成语言模型可读的嵌入，可参考\BookChapterRef{chap:vision-multimodal-models}{卷十的多模态模型章}。

训练目标与第\ref{chap:nlp-supervised-finetuning}章的监督微调相同，只对回答词元计算下一词元交叉熵（式\eqref{eq:nlp-sft-masked-loss}），音频嵌入和指令只作条件。训练通常分阶段进行。先在识别、音频描述等大规模“声音—文字”配对上训练连接部分，使音频嵌入对齐到语言模型能读懂的空间；再在问答、翻译、情绪识别等多任务指令数据上微调，使模型按指令处理声音。

Qwen-Audio与SALMONN是这一路线的两种实现\cite{nlp-speech-chu2023qwenaudio,nlp-speech-tang2024salmonn}。Qwen-Audio用Whisper-large-v2初始化音频编码器（32层，约6.4亿参数），其后接步长2的池化层，使每个输出向量对应约40\,ms声音，30\,s录音得到750个向量，再接Qwen-7B语言模型。它在三十多种音频任务上联合训练，为减轻不同数据集标注格式的相互干扰，在输出前加入分层任务标签：转写或分析标记、音频语言、任务类型、输出文字语言、是否输出时间戳，最后是具体的输出格式说明。共享的标签促使相近任务共享知识，不同的标签把格式冲突的数据区分开。多任务预训练阶段冻结语言模型、只训练音频编码器；随后的对话版本Qwen-Audio-Chat冻结编码器、只微调语言模型。

SALMONN同时接入两个编码器：Whisper-large-v2擅长语音，通用音频编码器BEATs擅长环境声与音乐。两者都输出50\,Hz特征，逐帧拼接后送入窗口级Q-Former，窗口长$L=17$帧（约0.33\,s），每窗$N=1$个查询，30\,s录音得到88个向量。语言模型是冻结的Vicuna-13B，只在其自注意力的查询与值矩阵上加秩为8的LoRA（第\ref{sec:nlp-sft-peft}节），可训练参数约3300万。训练分三阶段：先用识别与音频描述数据预训练，再用十余种任务的指令数据微调，最后做“激活微调”。第三阶段针对的问题是任务过拟合：模型在前两阶段见到的多为短而固定的转写与描述，遇到开放指令时仍倾向于输出转写，忽略指令本身。SALMONN在推断时降低LoRA缩放系数，得到较长、较多样的回答，再用少量这样生成的样本微调。

两种设计的差别落在降采样位置与训练分工上。池化保留每一帧附近的全部信息，向量数较多；窗口级Q-Former把每0.33\,s压成一个向量，序列短，但细粒度的发音信息可能丢失。Qwen-Audio训练编码器而冻结语言模型，适合大规模多任务对齐；SALMONN冻结两侧主体，只训练连接器与LoRA，成本低。半级联模型擅长理解：可以回答“说话人是否生气”“背景里是什么声音”这类转写无法回答的问题。它的输出仍是文字，声音回复需要外部合成或第\ref{sec:nlp-speech-tts}节的词元化合成；Whisper类编码器以整段为单位编码，也不适合边听边答。

\subsection{输入流、轮次与打断}
\label{sec:nlp-speech-turns}

无论哪条路线，系统都要判断三件不同的事：现在有没有人说话，用户的话轮是否结束，这句话起什么作用。“帮我订……明天下午”中的停顿可能是思考；系统播报时用户说“等等”，是要求接管或更正；短促的“嗯”则可能只表示在听。在线输入需要保留双方声音的时间轴、系统自身的播放状态与对话历史，监督数据也分别标注声音活动区间与轮次事件。下面先介绍只看已发生声音的检测，再介绍预测未来说话者活动的模型。

\subsubsection{VAD与轮次结束检测}
\label{sec:nlp-speech-vad}

\term{语音活动检测}（voice activity detection, VAD）逐帧判断是否存在语音。最简单的能量基线对第$i$个短窗计算$E_i=w^{-1}\sum_{n=0}^{w-1}x_{ih+n}^2$，其中$w$为窗长、$h$为帧移。$E_i$超过启动阈值$\theta_{\mathrm{on}}$时进入活动状态，连续若干窗低于结束阈值$\theta_{\mathrm{off}}$才退出；取$\theta_{\mathrm{on}}>\theta_{\mathrm{off}}$形成滞回，避免在阈值附近来回抖动。能量高的可能是噪声，能量低的可能是轻声，所以实际系统多用神经VAD：以有声、无声帧标注训练一个小型帧分类器，再做同样的滞回平滑。

轮次结束检测要回答的是话是否说完，它需要声音活动以外的证据。一个常用实现是“听取—候选结束—等待—确认结束”状态机：VAD转为静音时进入候选结束，再结合静音时长、句末韵律下降与语义完整度决定是否确认；若用户在等待期间恢复发声，就回到听取，撤销尚未提交的结束判断。教学设定阈值为400\,ms：用户在“帮我订”后停顿300\,ms再说“明天下午”，静音未达阈值，系统继续听；即使停顿达到500\,ms，“帮我订”缺少预约对象，语义完整度低，系统应发起澄清而不是执行。评价时分别统计过早截断（用户没说完就开始回答）与结束等待（用户说完到系统确认的时间），两者此消彼长，由阈值调节。系统自己的播放声音也会回到麦克风，回声消除或双通道录音决定VAD实际接收什么。

\subsubsection{未来声音活动预测：VAP}
\label{sec:nlp-speech-vap}

静音时长只是已经发生的证据，人在对话中却会提前预判对方何时说完。\term{声音活动预测}（voice activity projection, VAP）把轮次控制改写为预测：根据已见声音与双方活动历史，预测未来2\,s内两人各自的活动模式\scite{nlp-speech-ekstedt2022vap}。原模型把单声道混合录音送入预训练CPC编码器，得到100\,Hz的256维帧表示；再加上双方当前的活动标记与过去若干时间窗的活动比例，送入4层、宽256的因果Transformer。

预测目标是离散的联合状态。未来2\,s划为200、400、600、800\,ms四个由细到粗的区间，每人每个区间按活动比例是否超过阈值记为0或1，两人共8位，对应$2^8=256$个联合状态。模型输出256类分布，用交叉熵训练；标签由录音的后续活动自动构造，不需要人工标注，模型输入只含当前与过去，不会泄漏未来。与对8个区间分别做二分类相比，联合分布保留了两人活动之间的依赖，例如“A停、B接”与“A、B同时停”可以有不同的概率。

推断时，把与某个事件相符的状态概率相加，就得到该事件的概率，不需要为每种事件单独训练。在双方都静音的时刻，比较“只有A将说话”与“只有B将说话”两组状态的总概率，就得到保持话轮或移交话轮的判断；对“对方只在近期区间短暂活动”的状态求和，得到短回应的概率。教学例子：设符合移交、保持与其他模式的总概率为$0.55,0.35,0.10$，在“移交或保持”二选一的条件下，移交概率为$0.55/(0.55+0.35)\approx0.611$，而不加条件的移交概率仍是$0.55$；报告时须说明用的是哪一种。VAP给出的是何时可能轮到系统说话，它不判断预约信息是否齐全，后者要由下一小节的语义状态决定。

\subsection{增量理解与响应选择}
\label{sec:nlp-speech-incremental-understanding}

流式识别会把“明天”改成“后天”，转写不是只追加的日志。每条局部转写或语义记录都带版本号、对应的源声音区间，以及已确认前缀与暂定尾部。状态更新按新增、替换与撤回三种操作进行：版本1暂记4月18日；版本2听到“不，后天”后改为4月19日，依赖旧日期且尚未执行的动作同时失效。原话可以保留供追溯，当前有效值只有一份。在这一状态上，系统要在等待、短确认、澄清与正式回答之间选择响应。

\subsubsection{状态驱动的响应门控}
\label{sec:nlp-speech-response-gate}

响应门控按状态与事件选择下一个动作。规则实现依次检查五个条件：语义是否完整，关键字段是否确认，用户是否交出话轮（来自上一小节的轮次检测或VAP），工具是否返回，以及剩余延迟预算。算法\ref{alg:nlp-speech-response-gate}给出一个教学版本。规则来自业务前置条件与交互协议；有标注的历史对话时，也可以在同一动作集上训练分类器，或让语言模型提出动作，再经过同样的字段、权限与结果检查。

\begin{algorithm}[状态驱动的响应门控（教学版本）]
  \label{alg:nlp-speech-response-gate}
  \begin{algorithmic}[1]
    \Require 语义状态$\mathcal S$（字段、版本、确认标记），轮次信号$q\in\{\text{未结束},\text{结束}\}$，在途工具请求集合$R$，剩余等待预算$b$
    \Ensure 动作$a\in\{\text{继续听},\text{短确认},\text{澄清},\text{请求授权},\text{正式回答}\}$
    \If{有新版本替换了字段}\Comment{更正优先}
      \State 取消依赖旧版本且未执行的派生动作；迟到结果按版本号标为过期
    \EndIf
    \If{$q=\text{未结束}$}
      \State \Return 继续听
    \EndIf
    \If{必需字段缺失}
      \State \Return 澄清（询问缺失字段）
    \EndIf
    \If{涉及写操作且未获授权}
      \State \Return 请求授权
    \EndIf
    \If{$R$中有当前版本的请求尚未返回}
      \State \Return 若$b$即将耗尽则短确认（“正在查询”），否则继续等待
    \EndIf
    \State \Return 正式回答（以当前版本的已返回结果为依据）
  \end{algorithmic}
\end{algorithm}

对“明天……不，后天”，第一次只创建暂定日期，第二次替换日期并取消依赖旧版本的待发请求；已经发出的查询返回时先核对版本，不把旧结果写进新状态。\term{短应答}（backchannel）如“嗯，我在听”只表示收到，不当作业务许可。教学门控看到“日期已确认、人数缺失、轮次结束”时，输出澄清人数并进入等待；看到“字段完整、查询在途”时，至多短确认，不编造可预约时段。评价除动作标签的准确率外，还要统计实际的抢话次数、把暂定值提交给工具的次数，以及澄清后缺失信息是否真的补回。

\subsection{全双工语音对话：Moshi}
\label{sec:nlp-speech-moshi}

级联与半级联系统都按话轮交替工作，何时说、何时停取决于外部检测与门控。Moshi换了一种建模方式：用户与系统两路声音都离散成词元，同一个模型在共同的80\,ms时间网格上，每步同时读取用户的新声音并生成系统的下一帧声音\scite{defossez2024moshi}。说话、沉默、交叠和短回应都只是两路声音在时间轴上的不同模式，不再需要单独的轮次决策模块。实现这一点需要三个部件：低帧率且带语义的音频编解码器Mimi，能在每步生成多个码流的RQ-Transformer，以及让文字先于声音出现的内心独白。

第\ref{sec:nlp-speech-codec}小节的神经编解码器以残差向量量化（式\eqref{eq:nlp-speech-rvq}）把每帧压成若干级码本索引。语言模型直接建模这些码流时，帧率决定序列长度，所以Mimi把帧率压到12.5\,Hz。编码器是因果卷积网络，各层步长依次为4、5、6、8、2，总步长1920，把24\,kHz波形变为每80\,ms一帧的512维潜表示；量化前后各有一个8层因果Transformer，增强帧间建模。量化用8个码本、每个2048项，码率为$12.5\times8\times\log_2 2048=1100$\,bit/s，约1.1\,kbps。

Mimi的第一级码本经过\term{语义蒸馏}（semantic distillation）：用WavLM（第\ref{sec:nlp-speech-hubert}小节）提取同一段录音的表示，平均池化到12.5\,Hz，把第一级量化输出线性投影后与之计算余弦距离，作为附加损失。若第一级仍处在普通残差量化链的开头，后续各级要量化它留下的残差，语义约束会与声学重建相互牵制。Mimi因此拆分量化器：一个单独的向量量化器负责语义码，一个7级残差量化器并行负责声学细节，两者输出相加后送入解码器。第一级因而成为语义词元，其余7级是声学词元。训练只用对抗损失与特征匹配损失，不用频谱重建损失；原文报告这样客观指标下降，主观音质反而提高。

每个80\,ms步要生成17个子词元：1个系统文字词元，系统声音的8级码，用户声音的8级码。若把它们展平成一条序列，7B规模的语言模型每秒要跑$12.5\times17=212.5$步。RQ-Transformer借用残差量化图像生成中的做法，把时间与码本深度分给两个Transformer\cite{nlp-speech-lee2022rqtransformer}。记第$s$步的子词元为$V_{s,1},\ldots,V_{s,K}$，$K=17$。时间Transformer读取此前所有步，每步的输入是$K$个子词元嵌入之和，输出上下文向量
\begin{equation}
 \bm z_s=\mathrm{Tr}_{\mathrm{Temp}}(V_0,\ldots,V_{s-1})\in\mathbb R^{d}.
 \label{eq:nlp-speech-rq-temporal}
\end{equation}
深度Transformer在步内沿$k$自回归，输出第$k$个子词元的logit：
\[
 \bm l_{s,1}=\mathrm{Lin}(\bm z_s),\qquad
 \bm l_{s,k}=\mathrm{Tr}_{\mathrm{Depth}}(\bm z_s,V_{s,1},\ldots,V_{s,k-1}),\quad k>1,
\]
$\operatorname{softmax}(\bm l_{s,k})$近似$V_{s,k}$在全部历史与本步前$k-1$个子词元条件下的分布。联合分布于是分解为
\begin{equation}
 p(V)=\prod_{s}\prod_{k=1}^{K}p\bigl(V_{s,k}\mid V_{<s},V_{s,<k}\bigr).
 \label{eq:nlp-speech-rq-factorization}
\end{equation}
时间Transformer与7B文本模型Helium结构相同（32层、宽4096），并由Helium初始化，上下文取3000步（约4分钟），从而继承文本知识；深度Transformer只有6层、宽1024，各$k$使用独立的投影参数。大模型每80\,ms只运行一次，小模型在步内运行17次，计算量远小于展平序列。图\ref{fig:nlp-speech-moshi-rq}画出了一个时间步内的计算。

\begin{figure*}[htbp]
  \centering
  % 第九卷 15.7.4 节：Moshi 在一个 80 ms 时间步内的 RQ-Transformer 计算（通栏，约 145 mm）。
% 蓝色：用户侧输入；黄色描边：两个 Transformer；红色：系统生成的词元与声音。
% 左侧时间 Transformer 每步运行一次，右侧深度 Transformer 在步内依次预测 K=17 个子词元。
\begin{tikzpicture}[
  x=1mm,y=1mm,
  figure picture,
  model/.style={figure node,draw=FigureModel,fill=FigureModelFill,
    minimum height=12mm,font=\figurefont\scriptsize},
  inp/.style={figure node,draw=FigureInput,fill=FigureInputFill,
    minimum height=8mm,font=\figurefont\scriptsize},
  tok/.style={figure node,draw=FigureOutput,fill=FigureOutputFill,
    minimum height=7mm,minimum width=12mm,inner xsep=1.5pt,
    font=\figurefont\scriptsize},
  utok/.style={tok,draw=FigureInput,fill=FigureInputFill},
  lab/.style={font=\figurefont\scriptsize,text=FigureMuted,align=center},
  fl/.style={figure flow},
  back/.style={figure flow,draw=FigureMuted,dash pattern=on 2pt off 1.5pt}
]
  % 时间维：每 80 ms 运行一次的大模型。
  \node[inp,minimum width=36mm] (hist) at (20,0) {历史步$V_0,\ldots,V_{s-1}$\\每步17个词元};
  \node[model,minimum width=36mm] (emb) at (20,15) {每步17个嵌入相加\\得到$\mathbb R^{4096}$输入};
  \node[model,minimum width=36mm,minimum height=16mm] (temp) at (20,34)
    {时间Transformer\\Helium初始化\\32层，$d=4096$};
  \draw[fl] (hist) -- (emb);
  \draw[fl] (emb) -- (temp);

  % 深度维：步内依次预测 17 个子词元的小模型。
  \node[model,minimum width=84mm,minimum height=14mm] (depth) at (100,34)
    {深度Transformer（6层，$d=1024$，各$k$有独立投影）\\读取$\bm z_s$与本步已生成的$V_{s,1},\ldots,V_{s,k-1}$，输出$V_{s,k}$的分布};
  \draw[fl] (temp.east) -- node[figure edge label,above=0.5mm] {$\bm z_s$} (depth.west);

  % 第 s 步的 17 个子词元。
  \node[tok] (w) at (64,58) {$W_s$\\文字};
  \node[tok] (a1) at (79,58) {$A_{s,1}$\\语义};
  \node[tok,minimum width=19mm] (aq) at (97,58) {$A_{s-\tau,2:8}$\\声学7级};
  \node[utok] (u1) at (116,58) {$A'_{s,1}$\\用户语义};
  \node[utok,minimum width=19mm] (uq) at (134,58) {$A'_{s-\tau,2:8}$\\用户声学};
  \draw[fl] (w) -- (a1);
  \draw[fl] (a1) -- (aq);
  \draw[fl] (aq) -- (u1);
  \draw[fl] (u1) -- (uq);
  \foreach \n in {w,a1,aq,u1,uq} {\draw[fl] (depth.north -| \n) -- (\n.south);}
  \node[figure edge label,font=\figurefont\scriptsize] at (64,47) {$k=1$};
  \node[figure edge label,font=\figurefont\scriptsize] at (79,47) {$k=2$};
  \node[figure edge label,font=\figurefont\scriptsize] at (97,47) {$k=3,\ldots,9$};
  \node[figure edge label,font=\figurefont\scriptsize] at (125,47) {$k=10,\ldots,17$};

  % 输出与输入的声音。
  \node[figure node,draw=FigureOutput,fill=FigureOutputFill,minimum width=32mm,
    font=\figurefont\scriptsize] (dec) at (88,82) {Mimi解码器\\系统声音，80\,ms一帧};
  \draw[fl] (a1.north) -- (dec.south -| a1);
  \draw[fl] (aq.north) -- (dec.south -| aq);
  \node[inp,minimum width=36mm] (enc) at (125,82) {用户声音$\to$Mimi编码器\\推断时取实际输入};
  \draw[fl,draw=FigureInput] (enc.south -| u1) -- (u1.north);
  \draw[fl,draw=FigureInput] (enc.south -| uq) -- (uq.north);

  % 本步结果写回历史。
  \draw[back] (w.west) -- (50,58) |- (-4,72) -- (-4,0) -- (hist.west);
  \node[lab,anchor=south west] at (-3,72.5) {第$s$步的17个词元写回历史，供下一步读取};
\end{tikzpicture}

  \caption{Moshi在第$s$个80\,ms时间步内的计算。左侧时间Transformer读取此前所有步（每步17个子词元的嵌入之和），输出$\bm z_s$；右侧深度Transformer以$\bm z_s$为条件，按$k=1,\ldots,17$依次预测文字词元、系统语义与声学词元、用户语义与声学词元，声学词元相对语义词元延迟$\tau$步。系统词元经Mimi解码为声音；推断时用户词元取自实际收到的声音，不由模型采样。本步的17个词元写回历史，供下一步读取。}
  \label{fig:nlp-speech-moshi-rq}
\end{figure*}

只建模声音词元时，生成的语句常常语言质量差。\term{内心独白}（inner monologue）让系统在每步先生成与自身声音对齐的文字词元，作为声音的先导。训练数据用Whisper转写系统一侧的录音，得到带时间戳的词；第$i$个词的起始帧$t_i$取其起始时间乘以12.5，它的子词依次放在$t_i,t_i+1,\ldots$；其余步填PAD，词前一步放EPAD标记词的开始。英语对话中约65\%的文字位置是PAD。用户一侧不转写，文字流只描述系统自己要说的话。

声学码依赖语义码：同一帧的细节要等内容确定后才好预测。Moshi让声学码相对语义码延迟$\tau$步，第$s$步的子词元为
\begin{equation}
 V_{s}=\bigl(W_s,\;A_{s,1},\;A_{s-\tau,2},\ldots,A_{s-\tau,8},\;A'_{s,1},\;A'_{s-\tau,2},\ldots,A'_{s-\tau,8}\bigr),
 \label{eq:nlp-speech-moshi-delay}
\end{equation}
$W_s$是文字词元，$A_{s,q}$是系统声音第$q$级码，$A'_{s,q}$是用户声音的对应码，$s-\tau<1$时置为零。于是预测第$s-\tau$帧的声学码时，模型已经看到该帧之后$\tau$步的语义码。原文预训练取$\tau=2$，后续微调取$\tau=1$。文字流与声音之间的相对延迟还可以改变模型的用途：文字领先声音时模型相当于流式合成，声音领先文字时相当于流式识别。

训练目标是全部子词元的交叉熵。文字词元与全部声音词元各占一半权重，声音词元内部再按类型加权：
\[
 L=\frac1S\sum_{s=1}^{S}\Bigl(\operatorname{CE}(\bm l_{s,1},V_{s,1})+\frac{\sum_{k=2}^{K}\alpha_k\operatorname{CE}(\bm l_{s,k},V_{s,k})}{\sum_{k=2}^{K}\alpha_k}\Bigr),
\]
语义码取$\alpha_k=100$、声学码取$\alpha_k=1$，使有限的模型容量优先保证说出的内容。训练时用户一路的词元也计入损失；推断时只采样系统的文字与声音，用户一路直接填入Mimi对实际用户声音的编码。这样模型说话时用户流仍在输入，交叠与打断都表示在两路词元之中。

12.5\,Hz帧率规定了表示的节拍：零延迟时最小延迟为一帧80\,ms，每多一步声学延迟加80\,ms。原文据此给出$\tau=1$时160\,ms的理论延迟，实际约200\,ms。这个数字不含网络传输与播放缓冲，端到端延迟须按墙钟另测。对系统说“按明天……”时用户插话“不是，后天”的情形，要分别检查用户声音是否进入了模型，旧回复是否停下，后续声音是否改用正确日期，以及业务状态是否同步。

与级联系统相比，Moshi的优势在于自然的交叠、短回应与低延迟，语气与情绪在两侧都保留在声音词元中。代价是对话行为由模型自行决定，很难像算法\ref{alg:nlp-speech-response-gate}那样显式规定何时必须澄清或请求授权；文字流虽然可读，工具调用、事实核验与业务状态仍需外接，并通过下一小节的协调机制与声音输出对齐。

\subsection{声音回复与行为协调}
\label{sec:nlp-speech-coordination}

声音一旦播出就无法撤回。实时回复的每个片段因此按待生成、已生成未播、已播三种状态管理，并附上任务版本、依赖字段与播放位置。更正或新的工具结果使旧版本失效时，取消过时的生成任务，清空未播队列，并记录已经播出的范围；已播的错误内容只能明确更正。用户确实试图接管时停止或衰减当前播放，短应答则不触发停播，否则用户每说一次“嗯”，系统都会丢掉有用的回复。

工具请求同样附请求标识、任务版本与执行状态。取消本地等待不等于取消服务端的写操作；收到迟到结果时先查询实际执行记录，不直接重发写请求。幂等、结果查询与失败恢复由第\ref{chap:nlp-solving-and-execution}章的执行闭环处理，声音层保存这些证据并决定何时播报。

构造一条墙钟记录：用户在3.0\,s说完，系统3.8\,s开始播放“正在查询”，工具4.6\,s返回，结果声音5.0\,s开始播放。结束到首段声音为0.8\,s，结束到结果声音为2.0\,s；短确认缩短了反馈等待，业务完成仍在4.6\,s。再设用户在4.0\,s插话改日期，系统4.15\,s停止播放，打断停止延迟为150\,ms；4.6\,s返回的是旧日期的查询结果，只能存为过期版本，不能当作新日期的答案播出。

完整评价需要业务结果与交互质量同时计分。$\tau$-Voice在模拟业务环境中按环境终态检查任务是否完成，同时记录回应、让出话轮、打断与短回应处理等行为；它的模拟时间与墙钟解耦，适合组织故障测试，部署延迟仍要在真实墙钟下测量\scite{nlp-intro-ray2026}。实际测试应保存用户结束、稳定识别、首段播放、结果播放、打断停止与业务完成等事件，口音、噪声、交叠与接口失败都应进入完整交互测试；只有这些事件共享同一时间基准，表\ref{tab:nlp-speech-contracts}中的时间比较才成立。
